\section{Construction requirements and the remaining questions}\label{sec:open}
The results constrain two distinct resources: one-stage correctness dimension
and complete-state composition. Their combination removes several ways of
turning a large one-stage source into a complete algorithm, but it does not
supply an interface theorem for arbitrary native tape computations.

For \cite[Problem~9.1(a)]{SeventeenthPaper2026}, the next capacity question is whether every fixed number of shared deterministic
decoders has dimension at most $N/2+o(N)$ on linear sources. The two-class
$41/64$ bound leaves a constant gap; the fixed-class constructions give only
vanishing normalized excess. Nonlinear source promises and chosen-source
nonlinear-to-linear losses must not be substituted for this maximum-capacity
question. The measured covariance and symmetric-retention deficits provide
more structured intermediate targets. The exact arbitrary-set capacity shows
why a high-entropy fibre alone is insufficient: a native extraction needs the
linearity or actual affine-intersection premise used by its capacity bound.
The mask floors identify limits of the stated positive-term certificates;
improving beyond them requires more compatibility information, a smaller
actual starting ambient, or another argument.

The interval constructions attain the bound with deterministic binary
decoders when $d-M$ is even. This does not settle the original linear-decoder
attainment question for unequal pieces in
\cite[Problem~9.2]{SeventeenthPaper2026}. The diagonal-defect bound supplies
a way to retain actual block losses when testing that question.

For composition, a sharp class-product inequality for arbitrary linear masks
would improve the Ramsey cutoff. A positive normalized rank deficit would
strengthen the present growing rank loss. On the nonlinear side, the complete
masked-neighbor family remains open beyond the uniform, anchored and single-site
cases and the general period restrictions. The balance arguments apply to
actual conjugated operators and balanced fibres; an assigned covariance equation
or a surjective but unbalanced decoder is insufficient.

The complete scattered first factor is a concrete warning: total bijective
completion need not preserve a small class count in the required second map.
The residual bounds offer a different application route, if one can exhibit
a complete structured affine baseline and pay its joint correction record.
They do not supply such a baseline for an arbitrary fast computation.

Finally a native application must account for production of the representation,
access to its entries, the complete state at every intermediate stage, and
recovery of the ordered output. None of the shared-rule theorems proves that a
linear-time native transposer would possess the required fixed class count,
linear source or same-space factorization. This is the unresolved step toward
\cite[Problem~9.4]{SeventeenthPaper2026}. A native theorem giving one of those
interfaces with small paid advice could make the algebraic bounds decisive.
A construction must instead realize a complete compatible sequence of rules
within the native time budget. The two-stage theorems do not exclude a growing
number of stages, and none of the bounds supplies the multiscale time charge
asked for in \cite[Problem~9.5]{SeventeenthPaper2026}.
